\documentclass[12pt,a4paper]{article}

% ─── Packages ────────────────────────────────────────────────────────────────
\usepackage[utf8]{inputenc}
\usepackage[T1]{fontenc}
\usepackage[bahasa]{babel}
\usepackage{geometry}
\usepackage{graphicx}
\usepackage{hyperref}
\usepackage{listings}
\usepackage{xcolor}
\usepackage{booktabs}
\usepackage{array}
\usepackage{caption}
\usepackage{subcaption}
\usepackage{amsmath}
\usepackage{setspace}
\usepackage{fancyhdr}
\usepackage{titlesec}
\usepackage{enumitem}
\usepackage{float}
\usepackage{parskip}

% ─── Page Geometry ───────────────────────────────────────────────────────────
\geometry{
    top=3cm,
    bottom=3cm,
    left=4cm,
    right=3cm
}

% ─── Spacing ─────────────────────────────────────────────────────────────────
\onehalfspacing

% ─── Header / Footer ─────────────────────────────────────────────────────────
\pagestyle{fancy}
\fancyhf{}
\fancyhead[R]{\small Pengujian Perangkat Lunak -- Mockito}
\fancyfoot[C]{\thepage}
\renewcommand{\headrulewidth}{0.4pt}

% ─── Colors ──────────────────────────────────────────────────────────────────
\definecolor{codegreen}{rgb}{0.0,0.5,0.0}
\definecolor{codegray}{rgb}{0.5,0.5,0.5}
\definecolor{codepurple}{rgb}{0.58,0,0.82}
\definecolor{codeblue}{rgb}{0.0,0.0,0.6}
\definecolor{backcolour}{rgb}{0.97,0.97,0.97}
\definecolor{keywordcolor}{rgb}{0.0,0.0,0.6}
\definecolor{stringcolor}{rgb}{0.63,0.13,0.13}

% ─── Code Listing Style ──────────────────────────────────────────────────────
\lstdefinestyle{javaStyle}{
    backgroundcolor=\color{backcolour},
    commentstyle=\color{codegreen}\itshape,
    keywordstyle=\color{keywordcolor}\bfseries,
    numberstyle=\tiny\color{codegray},
    stringstyle=\color{stringcolor},
    basicstyle=\ttfamily\footnotesize,
    breakatwhitespace=false,
    breaklines=true,
    captionpos=b,
    keepspaces=true,
    numbers=left,
    numbersep=5pt,
    showspaces=false,
    showstringspaces=false,
    showtabs=false,
    tabsize=4,
    frame=single,
    rulecolor=\color{black!30},
    language=Java
}
\lstset{style=javaStyle}

% ─── Hyperref Setup ──────────────────────────────────────────────────────────
\hypersetup{
    colorlinks=true,
    linkcolor=black,
    urlcolor=blue,
    citecolor=black
}

% ─── Title Information ───────────────────────────────────────────────────────
\title{
    \vspace{-1cm}
    {\Large \textbf{LAPORAN PRAKTIKUM}} \\[0.5em]
    {\large \textbf{Pengujian Perangkat Lunak dengan Framework Mockito}} \\[0.3em]
    {\normalsize Studi Kasus: Sistem Manajemen Perpustakaan}
}
\author{
    \textbf{Program Studi Informatika} \\
    Pengujian Perangkat Lunak (PPPL)
}
\date{\today}

% ═════════════════════════════════════════════════════════════════════════════
\begin{document}
% ═════════════════════════════════════════════════════════════════════════════

\maketitle
\thispagestyle{fancy}

\begin{abstract}
\noindent
Laporan ini memaparkan hasil dan pembahasan pengujian perangkat lunak menggunakan
\textit{framework} Mockito pada studi kasus sistem manajemen perpustakaan yang
diimplementasikan dalam bahasa pemrograman Java. Pengujian dilakukan dengan
memanfaatkan konsep \textit{mock object} dan \textit{stub} guna mengisolasi
unit-unit kode dari ketergantungan eksternal, sehingga pengujian dapat berjalan
secara mandiri, deterministik, dan \textit{reproducible}. Tiga skenario utama
dikaji secara mendalam: (1) pengecekan ketersediaan buku berdasarkan
\textit{identifier} unik, (2) peminjaman buku dengan dua skenario kondisional,
dan (3) pengiriman notifikasi pasca peminjaman beserta verifikasi urutan
pemanggilan menggunakan \texttt{InOrder}. Seluruh pengujian dilaksanakan dengan
JUnit 5 dan Mockito 5.11.0 di atas JDK 17. Hasil pengujian menunjukkan bahwa
kedelapan \textit{test case} yang dirancang berhasil lulus (\textit{passed})
tanpa kegagalan, membuktikan kebenaran implementasi dan efektivitas pendekatan
\textit{test-driven} berbasis Mockito.
\end{abstract}

\tableofcontents
\newpage

% ─────────────────────────────────────────────────────────────────────────────
\section{Pendahuluan}
% ─────────────────────────────────────────────────────────────────────────────

\subsection{Latar Belakang}

Pengujian perangkat lunak (\textit{software testing}) merupakan salah satu
tahapan kritis dalam siklus hidup pengembangan perangkat lunak (SDLC) yang
bertujuan memastikan kualitas, keandalan, dan kebenaran fungsional suatu sistem.
Dalam konteks \textit{unit testing}, pengujian idealnya bersifat terisolasi:
setiap unit fungsi atau metode diuji secara independen tanpa bergantung pada
komponen lain yang mungkin belum selesai diimplementasikan atau memiliki
ketidakpastian dalam perilakunya.

Untuk mencapai isolasi tersebut, digunakanlah teknik \textit{mocking}, yaitu
pembuatan objek tiruan (\textit{mock object}) yang menggantikan peran objek
nyata selama pengujian berlangsung. Mockito merupakan \textit{framework}
\textit{mocking} yang paling banyak digunakan dalam ekosistem Java, menyediakan
API yang ekspresif dan mudah dipahami untuk mendefinisikan perilaku objek tiruan
(\textit{stubbing}) serta memverifikasi interaksi antar komponen
(\textit{mocking}/\textit{verification}).

\subsection{Tujuan}

Tujuan dari laporan ini adalah:
\begin{enumerate}[label=\arabic*.]
    \item Menjelaskan secara akademis konsep \textit{stub} dan \textit{mock}
          dalam konteks pengujian perangkat lunak berbasis Mockito.
    \item Memaparkan implementasi dan hasil pengujian fungsi pengecekan
          ketersediaan buku pada lapisan \textit{controller} dan \textit{model}.
    \item Menganalisis dua skenario pengujian peminjaman buku:
          kondisi buku tersedia dan kondisi buku sedang dipinjam.
    \item Membahas penggunaan \texttt{InOrder} Mockito untuk memverifikasi
          urutan pemanggilan metode dalam alur peminjaman buku beserta
          pengiriman notifikasi.
\end{enumerate}

\subsection{Ruang Lingkup}

Studi kasus yang digunakan adalah sistem manajemen perpustakaan sederhana dengan
tiga kelas utama: \texttt{LibraryModel} (lapisan data/akses data),
\texttt{LibraryControl} (lapisan logika bisnis/\textit{controller}), dan
\texttt{NotificationService} (layanan pengiriman notifikasi). Pengujian
dilaksanakan pada kelas \texttt{LibraryControl} dengan dependensinya
(\texttt{LibraryModel} dan \texttt{NotificationService}) diejek menggunakan
Mockito.

% ─────────────────────────────────────────────────────────────────────────────
\section{Tinjauan Pustaka}
% ─────────────────────────────────────────────────────────────────────────────

\subsection{Unit Testing}

\textit{Unit testing} adalah proses verifikasi bahwa unit terkecil dari
perangkat lunak—umumnya sebuah metode atau fungsi—bekerja sesuai dengan
spesifikasi yang ditetapkan. Ciri khas \textit{unit test} yang baik mencakup:
kecepatan eksekusi tinggi, isolasi penuh dari komponen lain (\textit{FIRST
principle: Fast, Independent, Repeatable, Self-validating, Timely}), serta
kemudahan pemeliharaan.

\subsection{Mock Object dan Stub}

Dalam taksonomi \textit{test doubles} yang diperkenalkan oleh Gerard Meszaros,
terdapat beberapa jenis pengganti objek nyata, dua di antaranya yang paling
relevan adalah:

\begin{itemize}
    \item \textbf{Stub}: Objek tiruan yang menyediakan respons telah ditentukan
          sebelumnya (\textit{canned responses}) terhadap pemanggilan metode
          tertentu selama pengujian. Stub tidak memverifikasi bagaimana metode
          dipanggil; ia hanya memastikan bahwa \textit{system under test} (SUT)
          menerima data yang dibutuhkan untuk menyelesaikan eksekusinya.
          Dalam Mockito, \textit{stubbing} dilakukan menggunakan sintaks
          \texttt{Mockito.when(\ldots).thenReturn(\ldots)}.

    \item \textbf{Mock}: Objek tiruan yang diprogram dengan ekspektasi—berisi
          definisi tentang panggilan metode mana yang \emph{harus} terjadi
          dan \emph{berapa kali}. Setelah SUT dieksekusi, mock dapat
          diverifikasi untuk memastikan interaksi yang diharapkan benar-benar
          terjadi. Dalam Mockito, verifikasi dilakukan menggunakan
          \texttt{Mockito.verify(\ldots)}.
\end{itemize}

\subsection{Framework Mockito}

Mockito adalah \textit{open-source mocking framework} untuk Java yang
memungkinkan pembuatan dan konfigurasi objek tiruan secara deklaratif.
Fitur-fitur utama Mockito meliputi:

\begin{itemize}
    \item Anotasi \texttt{@Mock} untuk deklarasi \textit{mock object}.
    \item Anotasi \texttt{@InjectMocks} untuk injeksi otomatis \textit{mock}
          ke dalam SUT.
    \item \texttt{Mockito.when().thenReturn()} untuk \textit{stubbing}.
    \item \texttt{Mockito.verify()} dengan berbagai mode verifikasi seperti
          \texttt{times(n)}, \texttt{never()}, \texttt{atLeast(n)}.
    \item \texttt{InOrder} untuk verifikasi urutan pemanggilan metode.
    \item \texttt{ArgumentMatchers} seperti \texttt{anyInt()},
          \texttt{anyString()} untuk \textit{flexible argument matching}.
\end{itemize}

% ─────────────────────────────────────────────────────────────────────────────
\section{Arsitektur dan Implementasi Sistem}
% ─────────────────────────────────────────────────────────────────────────────

\subsection{Struktur Proyek}

Proyek dibangun menggunakan Apache Maven dengan konfigurasi kompilator Java 17
dan dependensi JUnit Jupiter 5.10.2 serta Mockito 5.11.0. Struktur direktori
proyek adalah sebagai berikut:

\begin{verbatim}
PPPL_Mockito/
├── pom.xml
└── src/
    ├── main/java/org/example/
    │   ├── Book.java
    │   ├── LibraryModel.java
    │   ├── LibraryControl.java
    │   └── NotificationService.java
    └── test/java/
        └── LibraryTest.java
\end{verbatim}

\subsection{Kelas \texttt{Book}}

Kelas \texttt{Book} merepresentasikan entitas buku dalam sistem perpustakaan.
Atribut yang dimiliki mencakup \texttt{id} (identifikasi unik buku),
\texttt{title} (judul buku), \texttt{author} (nama pengarang), dan
\texttt{borrowed} (status peminjaman berupa nilai boolean).

\begin{lstlisting}[caption={Implementasi kelas \texttt{Book}},
                   label={lst:book}]
package org.example;

public class Book {
    private int id;
    private String title;
    private String author;
    private boolean borrowed;

    public Book(int id, String title, String author) {
        this.id = id;
        this.title = title;
        this.author = author;
        this.borrowed = false;
    }

    public int getId()            { return id; }
    public String getTitle()      { return title; }
    public String getAuthor()     { return author; }
    public boolean isBorrowed()   { return borrowed; }
    public void setBorrowed(boolean borrowed) {
        this.borrowed = borrowed;
    }
}
\end{lstlisting}

\subsection{Kelas \texttt{LibraryModel}}

\texttt{LibraryModel} berperan sebagai lapisan akses data (\textit{data access
layer}). Kelas ini bertanggung jawab atas operasi-operasi yang berhubungan
langsung dengan penyimpanan dan pengambilan data buku. Dalam konteks pengujian,
\texttt{LibraryModel} menjadi objek yang diejek (\textit{mocked}) karena
bergantung pada mekanisme persistensi data yang berada di luar cakupan
\textit{unit testing}.

\begin{lstlisting}[caption={Implementasi kelas \texttt{LibraryModel}},
                   label={lst:model}]
package org.example;

import java.util.ArrayList;
import java.util.List;

public class LibraryModel {
    public List<Book> getAllBooks() {
        return new ArrayList<>();
    }

    public void saveSearchKeyword(String keyword) { }

    public boolean checkAvailability(int id) {
        return false;
    }

    public void borrowBook(int id) { }
}
\end{lstlisting}

\subsection{Kelas \texttt{LibraryControl}}

\texttt{LibraryControl} merupakan lapisan logika bisnis (\textit{business logic
layer}) yang mengorkestrasikan interaksi antara antarmuka pengguna dan lapisan
data. Kelas ini menerima dependensi \texttt{LibraryModel} dan
\texttt{NotificationService} melalui mekanisme \textit{constructor injection},
mematuhi prinsip \textit{Dependency Inversion} dari SOLID.

\begin{lstlisting}[caption={Implementasi kelas \texttt{LibraryControl}},
                   label={lst:control}]
package org.example;

import java.util.List;

public class LibraryControl {
    private final LibraryModel model;
    private NotificationService notificationService;

    public LibraryControl(LibraryModel model) {
        this.model = model;
    }

    public LibraryControl(LibraryModel model,
                          NotificationService notificationService) {
        this.model = model;
        this.notificationService = notificationService;
    }

    public String searchBookAuthor(String title) {
        List<Book> books = model.getAllBooks();
        String result = "";
        for (Book book : books) {
            if (book.getTitle().equals(title)) {
                model.saveSearchKeyword(title);
                result = book.getAuthor();
            }
        }
        return result;
    }

    public boolean checkAvailability(int id) {
        return model.checkAvailability(id);
    }

    public boolean borrowBook(int id) {
        boolean available = model.checkAvailability(id);
        if (available) {
            model.borrowBook(id);
            return true;
        }
        return false;
    }

    public boolean borrowBookWithNotification(String title,
                                              String username) {
        List<Book> books = model.getAllBooks();
        for (Book book : books) {
            if (book.getTitle().equals(title) && !book.isBorrowed()) {
                model.borrowBook(book.getId());
                notificationService.sendNotification(username, title);
                return true;
            }
        }
        return false;
    }
}
\end{lstlisting}

\subsection{Kelas \texttt{NotificationService}}

\texttt{NotificationService} merupakan layanan yang bertugas mengirimkan
notifikasi kepada pengguna setelah transaksi peminjaman buku berhasil
dilakukan. Dalam lingkungan produksi, kelas ini dapat terhubung ke sistem
email, SMS, atau \textit{push notification}. Pada konteks pengujian, kelas ini
diejek untuk menghindari pengiriman notifikasi nyata dan untuk memverifikasi
bahwa metode \texttt{sendNotification} dipanggil dengan argumen yang benar.

\begin{lstlisting}[caption={Implementasi kelas \texttt{NotificationService}},
                   label={lst:notification}]
package org.example;

public class NotificationService {
    public void sendNotification(String username, String title) {
        System.out.println("Notification sent to " + username
            + ": Buku '" + title
            + "' berhasil dipinjam/dikembalikan.");
    }
}
\end{lstlisting}

% ─────────────────────────────────────────────────────────────────────────────
\section{Hasil dan Pembahasan}
% ─────────────────────────────────────────────────────────────────────────────

Subbagian berikut membahas tiga skenario utama pengujian yang dirancang sesuai
dengan kebutuhan sistem manajemen perpustakaan. Setiap skenario dianalisis dari
perspektif identifikasi \textit{stub}/\textit{mock}, logika implementasi, serta
interpretasi hasil pengujian.

% ─── 4.1 ─────────────────────────────────────────────────────────────────────
\subsection{Skenario 1: Pengecekan Ketersediaan Buku}

\subsubsection{Deskripsi Fungsional}

Skenario ini menguji kemampuan sistem untuk menentukan apakah sebuah buku
dengan \textit{identifier} tertentu (\texttt{id}) tersedia untuk dipinjam atau
tidak. Metode \texttt{checkAvailability(int id)} pada \texttt{LibraryControl}
mendelegasikan pengecekan kepada \texttt{LibraryModel}, yang berinteraksi
langsung dengan mekanisme persistensi data.

\subsubsection{Identifikasi Stub dan Mock}

\begin{table}[H]
\centering
\caption{Identifikasi \textit{Test Double} pada Skenario Pengecekan Ketersediaan}
\label{tab:stub-mock-availability}
\begin{tabular}{@{}p{3.5cm}p{2.5cm}p{4.5cm}p{3cm}@{}}
\toprule
\textbf{Objek} & \textbf{Jenis} & \textbf{Konfigurasi} & \textbf{Peran} \\
\midrule
\texttt{bookModel} &
    Stub \& Mock &
    \texttt{when(bookModel.checkAvailability(1)).thenReturn(true)} &
    Menyediakan nilai kembalian palsu sekaligus diverifikasi interaksinya \\
\addlinespace
\texttt{bookModel} &
    Mock &
    \texttt{verify(bookModel, times(1)).checkAvailability(1)} &
    Memastikan metode dipanggil tepat satu kali \\
\bottomrule
\end{tabular}
\end{table}

\subsubsection{Implementasi Test Method}

Dua \textit{test method} dirancang untuk mencakup kondisi positif (buku
tersedia) dan kondisi negatif (buku tidak tersedia):

\begin{lstlisting}[caption={Test method pengecekan ketersediaan buku},
                   label={lst:test-availability}]
@Test
void testCheckAvailability_bookAvailable() {
    // Stub: mendefinisikan perilaku model saat dipanggil dengan id=1
    Mockito.when(bookModel.checkAvailability(1)).thenReturn(true);

    boolean result = control.checkAvailability(1);

    // Assertion: memverifikasi nilai kembalian
    Assertions.assertTrue(result);
    // Mock verification: memastikan metode dipanggil 1 kali
    Mockito.verify(bookModel, Mockito.times(1)).checkAvailability(1);
}

@Test
void testCheckAvailability_bookNotAvailable() {
    // Stub: id=99 merepresentasikan buku yang tidak ditemukan/tersedia
    Mockito.when(bookModel.checkAvailability(99)).thenReturn(false);

    boolean result = control.checkAvailability(99);

    Assertions.assertFalse(result);
    Mockito.verify(bookModel, Mockito.times(1)).checkAvailability(99);
}
\end{lstlisting}

\subsubsection{Analisis Hasil}

Pada \texttt{testCheckAvailability\_bookAvailable}, \textit{stubbing}
\texttt{when(bookModel.checkAvailability(1)).thenReturn(true)} mengkonfigurasi
objek mock \texttt{bookModel} agar mengembalikan nilai \texttt{true} ketika
dipanggil dengan argumen \texttt{id = 1}. Dengan demikian, metode
\texttt{control.checkAvailability(1)} yang secara internal memanggil
\texttt{model.checkAvailability(id)} akan menerima nilai \texttt{true} dan
meneruskannya kepada pemanggil. \textit{Assertion} \texttt{assertTrue(result)}
dikonfirmasi benar, dan verifikasi mock mengonfirmasi bahwa
\texttt{checkAvailability(1)} dipanggil tepat satu kali.

Pada \texttt{testCheckAvailability\_bookNotAvailable}, skenario yang diuji
adalah ketika buku dengan \texttt{id = 99} tidak tersedia. \textit{Stub}
mengembalikan \texttt{false}, sehingga hasil \textit{assertion}
\texttt{assertFalse} terkonfirmasi. Kedua \textit{test case} ini secara efektif
memvalidasi perilaku \textit{boundary} dari metode \texttt{checkAvailability}.

% ─── 4.2 ─────────────────────────────────────────────────────────────────────
\subsection{Skenario 2: Peminjaman Buku}

\subsubsection{Deskripsi Fungsional}

Metode \texttt{borrowBook(int id)} pada \texttt{LibraryControl} mengimplementasikan
alur peminjaman buku yang terdiri dari dua langkah: (1) memeriksa ketersediaan
buku melalui \texttt{model.checkAvailability(id)}, dan (2) memanggil
\texttt{model.borrowBook(id)} apabila buku tersedia. Jika buku sedang dipinjam,
metode \texttt{borrowBook} pada model \emph{tidak boleh} dipanggil.

\subsubsection{Identifikasi Stub dan Mock}

\begin{table}[H]
\centering
\caption{Identifikasi \textit{Test Double} pada Skenario Peminjaman Buku}
\label{tab:stub-mock-borrow}
\begin{tabular}{@{}p{3.5cm}p{2.5cm}p{4.5cm}p{3cm}@{}}
\toprule
\textbf{Objek} & \textbf{Jenis} & \textbf{Konfigurasi} & \textbf{Peran} \\
\midrule
\texttt{bookModel} &
    Stub &
    \texttt{when(checkAvailability(1)).thenReturn(true)} &
    Mensimulasikan buku tersedia \\
\addlinespace
\texttt{bookModel} &
    Mock &
    \texttt{verify(bookModel, times(1)).borrowBook(1)} &
    Verifikasi peminjaman dipanggil \\
\addlinespace
\texttt{bookModel} &
    Stub &
    \texttt{when(checkAvailability(1)).thenReturn(false)} &
    Mensimulasikan buku tidak tersedia \\
\addlinespace
\texttt{bookModel} &
    Mock &
    \texttt{verify(bookModel, never()).borrowBook(anyInt())} &
    Verifikasi peminjaman tidak dipanggil \\
\bottomrule
\end{tabular}
\end{table}

\subsubsection{Implementasi Test Method}

\begin{lstlisting}[caption={Test method peminjaman buku -- dua skenario},
                   label={lst:test-borrow}]
// Skenario 2a: Buku tidak sedang dipinjam
@Test
void testBorrowBook_bookNotBorrowed_shouldCallBorrowBook() {
    // Stub: buku dengan id=1 tersedia
    Mockito.when(bookModel.checkAvailability(1)).thenReturn(true);

    boolean result = control.borrowBook(1);

    // Assertion: peminjaman berhasil
    Assertions.assertTrue(result);
    // Mock verification: borrowBook HARUS dipanggil 1 kali
    Mockito.verify(bookModel, Mockito.times(1)).borrowBook(1);
}

// Skenario 2b: Buku sedang dipinjam
@Test
void testBorrowBook_bookAlreadyBorrowed_shouldNotCallBorrowBook() {
    // Stub: buku dengan id=1 tidak tersedia (sedang dipinjam)
    Mockito.when(bookModel.checkAvailability(1)).thenReturn(false);

    boolean result = control.borrowBook(1);

    // Assertion: peminjaman gagal
    Assertions.assertFalse(result);
    // Mock verification: borrowBook TIDAK BOLEH dipanggil
    Mockito.verify(bookModel, Mockito.never())
           .borrowBook(Mockito.anyInt());
}
\end{lstlisting}

\subsubsection{Analisis Hasil}

\paragraph{Skenario 2a (Buku Tersedia).}
Ketika \texttt{checkAvailability(1)} di-\textit{stub} mengembalikan
\texttt{true}, alur eksekusi \texttt{borrowBook} akan masuk ke blok kondisi
\texttt{if (available)} dan memanggil \texttt{model.borrowBook(id)}.
Verifikasi \texttt{verify(bookModel, times(1)).borrowBook(1)} memastikan bahwa
pemanggilan ini benar-benar terjadi tepat satu kali. Metode mengembalikan
\texttt{true}, yang dikonfirmasi oleh \textit{assertion} \texttt{assertTrue}.

\paragraph{Skenario 2b (Buku Sedang Dipinjam).}
Ketika \texttt{checkAvailability(1)} mengembalikan \texttt{false}, alur
eksekusi tidak masuk ke blok kondisi, sehingga \texttt{model.borrowBook(id)}
tidak pernah dipanggil. Verifikasi \texttt{verify(bookModel, never()).borrowBook(anyInt())}
menggunakan \texttt{never()} untuk membuktikan bahwa metode tersebut memang
tidak dipanggil sama sekali. Penggunaan \texttt{anyInt()} sebagai
\textit{argument matcher} memastikan verifikasi berlaku untuk semua nilai
argumen integer, sehingga lebih komprehensif. Metode mengembalikan
\texttt{false}, yang dikonfirmasi oleh \textit{assertion} \texttt{assertFalse}.

Kedua skenario ini secara kolektif membuktikan kebenaran logika kondisional
dalam metode \texttt{borrowBook}, menghindarkan pemanggilan yang tidak
semestinya terhadap lapisan model.

% ─── 4.3 ─────────────────────────────────────────────────────────────────────
\subsection{Skenario 3: Peminjaman Buku dengan Notifikasi dan Verifikasi Urutan}

\subsubsection{Deskripsi Fungsional}

Skenario ini menguji metode \texttt{borrowBookWithNotification(String title,
String username)} yang mengintegrasikan alur peminjaman buku dengan pengiriman
notifikasi. Urutan pemanggilan yang diharapkan adalah:
\begin{enumerate}
    \item \texttt{model.getAllBooks()} -- pengambilan seluruh koleksi buku
    \item \texttt{model.borrowBook(id)} -- pencatatan peminjaman pada lapisan data
    \item \texttt{notificationService.sendNotification(username, title)} --
          pengiriman notifikasi kepada peminjam
\end{enumerate}

Urutan ini penting karena memastikan integritas data: notifikasi hanya
dikirimkan setelah data peminjaman berhasil dicatat.

\subsubsection{Identifikasi Stub dan Mock}

\begin{table}[H]
\centering
\caption{Identifikasi \textit{Test Double} pada Skenario Peminjaman dengan Notifikasi}
\label{tab:stub-mock-notification}
\begin{tabular}{@{}p{3.5cm}p{2.5cm}p{4.5cm}p{3cm}@{}}
\toprule
\textbf{Objek} & \textbf{Jenis} & \textbf{Konfigurasi} & \textbf{Peran} \\
\midrule
\texttt{bookModel} &
    Stub &
    \texttt{when(getAllBooks()).thenReturn(books)} &
    Menyediakan daftar buku fiktif \\
\addlinespace
\texttt{bookModel} &
    Mock (\textit{InOrder}) &
    \texttt{inOrder.verify(bookModel).getAllBooks()} &
    Verifikasi urutan pemanggilan \\
\addlinespace
\texttt{bookModel} &
    Mock (\textit{InOrder}) &
    \texttt{inOrder.verify(bookModel).borrowBook(1)} &
    Verifikasi peminjaman dalam urutan \\
\addlinespace
\texttt{notificationService} &
    Mock (\textit{InOrder}) &
    \texttt{inOrder.verify(notificationService).sendNotification(...)} &
    Verifikasi notifikasi dalam urutan \\
\bottomrule
\end{tabular}
\end{table}

\subsubsection{Implementasi Test Method}

\begin{lstlisting}[caption={Test method peminjaman dengan notifikasi dan InOrder},
                   label={lst:test-notification}]
@Test
void testBorrowBookWithNotification_inOrder() {
    // Persiapan: buat daftar buku dengan status belum dipinjam
    List<Book> books = new ArrayList<>();
    Book book = new Book(1, "Tutorial Figma", "Budi");
    book.setBorrowed(false);
    books.add(book);

    // Stub: getAllBooks() mengembalikan daftar buku yang telah disiapkan
    Mockito.when(bookModel.getAllBooks()).thenReturn(books);

    boolean result = control.borrowBookWithNotification(
                         "Tutorial Figma", "Andi");

    Assertions.assertTrue(result);

    // InOrder: mendefinisikan objek yang akan diverifikasi urutannya
    InOrder inOrder = Mockito.inOrder(bookModel, notificationService);

    // Verifikasi urutan: 1. getAllBooks, 2. borrowBook, 3. sendNotification
    inOrder.verify(bookModel).getAllBooks();
    inOrder.verify(bookModel).borrowBook(1);
    inOrder.verify(notificationService)
           .sendNotification("Andi", "Tutorial Figma");
}

@Test
void testBorrowBookWithNotification_bookAlreadyBorrowed_inOrder() {
    List<Book> books = new ArrayList<>();
    Book book = new Book(1, "Tutorial Figma", "Budi");
    book.setBorrowed(true);  // buku sedang dipinjam
    books.add(book);

    Mockito.when(bookModel.getAllBooks()).thenReturn(books);

    boolean result = control.borrowBookWithNotification(
                         "Tutorial Figma", "Andi");

    Assertions.assertFalse(result);

    // getAllBooks harus dipanggil sekali
    Mockito.verify(bookModel, Mockito.times(1)).getAllBooks();
    // borrowBook dan sendNotification TIDAK BOLEH dipanggil
    Mockito.verify(bookModel, Mockito.never())
           .borrowBook(Mockito.anyInt());
    Mockito.verify(notificationService, Mockito.never())
           .sendNotification(Mockito.anyString(), Mockito.anyString());
}
\end{lstlisting}

\subsubsection{Analisis Penggunaan InOrder}

Penggunaan \texttt{InOrder} dalam Mockito memungkinkan verifikasi bahwa
serangkaian pemanggilan metode terjadi dalam urutan yang presisi. Hal ini
berbeda dengan verifikasi individual menggunakan \texttt{verify}, yang hanya
memastikan sebuah metode dipanggil tanpa memperhatikan relasi temporalnya
terhadap metode lain.

Secara implementatif, objek \texttt{InOrder} dibuat melalui
\texttt{Mockito.inOrder(bookModel, notificationService)}, di mana argumen yang
diberikan adalah daftar mock yang urutan interaksinya ingin diverifikasi.
Pemanggilan \texttt{inOrder.verify(mock).method()} kemudian dilakukan sesuai
urutan yang diharapkan.

Apabila urutan aktual tidak sesuai—misalnya \texttt{sendNotification} terpanggil
sebelum \texttt{borrowBook}—maka Mockito akan melempar
\texttt{VerificationInOrderFailure}, menyebabkan \textit{test} gagal.
Mekanisme ini sangat berguna untuk memvalidasi invarian urutan yang kritis
secara semantis, seperti memastikan notifikasi hanya dikirim setelah data
persisten berhasil diperbarui.

\subsubsection{Analisis Hasil}

\paragraph{Skenario 3a (Buku Tersedia, Peminjaman Berhasil).}
Dengan \textit{stub} \texttt{getAllBooks()} mengembalikan daftar berisi satu
buku dengan status \texttt{borrowed = false}, metode
\texttt{borrowBookWithNotification} menemukan kecocokan berdasarkan judul,
memanggil \texttt{model.borrowBook(1)}, lalu
\texttt{notificationService.sendNotification("Andi", "Tutorial Figma")}.
Verifikasi \texttt{InOrder} berhasil mengonfirmasi ketiga pemanggilan terjadi
dalam urutan yang benar. Metode mengembalikan \texttt{true}.

\paragraph{Skenario 3b (Buku Sedang Dipinjam).}
Dengan status buku \texttt{borrowed = true}, kondisi
\texttt{!book.isBorrowed()} bernilai \texttt{false}, sehingga blok peminjaman
tidak dieksekusi. Hanya \texttt{getAllBooks()} yang dipanggil. Verifikasi
\texttt{never()} pada \texttt{borrowBook} dan \texttt{sendNotification}
mengonfirmasi bahwa kedua metode tersebut tidak dieksekusi, menjamin integritas
logika bisnis.

% ─────────────────────────────────────────────────────────────────────────────
\section{Ringkasan Hasil Pengujian}
% ─────────────────────────────────────────────────────────────────────────────

Seluruh \textit{test case} yang diimplementasikan dalam kelas \texttt{LibraryTest}
dieksekusi menggunakan Maven (\texttt{mvn test}) dan menghasilkan luaran sebagai
berikut:

\begin{verbatim}
[INFO] Running LibraryTest
[INFO] Tests run: 8, Failures: 0, Errors: 0, Skipped: 0
[INFO] BUILD SUCCESS
\end{verbatim}

Tabel \ref{tab:test-summary} merangkum seluruh \textit{test case} beserta
statusnya:

\begin{table}[H]
\centering
\caption{Ringkasan Hasil Pengujian}
\label{tab:test-summary}
\begin{tabular}{@{}clll@{}}
\toprule
\textbf{No.} & \textbf{Nama Test Method} & \textbf{Skenario} & \textbf{Status} \\
\midrule
1 & \texttt{TestSearchBookAuthor} &
    Pencarian penulis buku & \textcolor{codegreen}{\textbf{PASSED}} \\
2 & \texttt{testSearchBookAuthor\_withNever} &
    Pencarian buku tidak ditemukan & \textcolor{codegreen}{\textbf{PASSED}} \\
3 & \texttt{testCheckAvailability\_bookAvailable} &
    Buku tersedia & \textcolor{codegreen}{\textbf{PASSED}} \\
4 & \texttt{testCheckAvailability\_bookNotAvailable} &
    Buku tidak tersedia & \textcolor{codegreen}{\textbf{PASSED}} \\
5 & \texttt{testBorrowBook\_bookNotBorrowed\_...} &
    Peminjaman berhasil & \textcolor{codegreen}{\textbf{PASSED}} \\
6 & \texttt{testBorrowBook\_bookAlreadyBorrowed\_...} &
    Peminjaman ditolak & \textcolor{codegreen}{\textbf{PASSED}} \\
7 & \texttt{testBorrowBookWithNotification\_inOrder} &
    Peminjaman + notifikasi berurutan & \textcolor{codegreen}{\textbf{PASSED}} \\
8 & \texttt{testBorrowBookWithNotification\_...\_inOrder} &
    Buku dipinjam, tidak ada notifikasi & \textcolor{codegreen}{\textbf{PASSED}} \\
\bottomrule
\end{tabular}
\end{table}

% ─────────────────────────────────────────────────────────────────────────────
\section{Kesimpulan}
% ─────────────────────────────────────────────────────────────────────────────

Berdasarkan hasil implementasi dan pengujian yang telah dilakukan, dapat
ditarik beberapa kesimpulan sebagai berikut:

\begin{enumerate}[label=\arabic*.]
    \item \textbf{Efektivitas Mockito dalam Unit Testing.} Framework Mockito
          terbukti efektif dalam mengisolasi \textit{unit under test}
          (\texttt{LibraryControl}) dari dependensi eksternalnya
          (\texttt{LibraryModel} dan \texttt{NotificationService}),
          memungkinkan pengujian yang deterministik dan berfokus pada logika
          bisnis.

    \item \textbf{Peran Stub dan Mock.} \textit{Stub} (melalui
          \texttt{when().thenReturn()}) berperan dalam menyediakan data
          masukan yang terkendali kepada SUT, sedangkan \textit{mock}
          (melalui \texttt{verify()}) berperan dalam memvalidasi bahwa
          interaksi yang diharapkan antar komponen benar-benar terjadi.
          Kombinasi keduanya menghasilkan pengujian yang komprehensif.

    \item \textbf{Keunggulan InOrder untuk Urutan Kritis.} Penggunaan
          \texttt{InOrder} Mockito memberikan kemampuan tambahan untuk
          memverifikasi urutan temporal pemanggilan metode, yang sangat
          relevan dalam skenario di mana urutan operasi memiliki implikasi
          semantis penting—seperti memastikan notifikasi hanya dikirim
          setelah transaksi data berhasil.

    \item \textbf{Kualitas Kode Uji.} Seluruh kedelapan \textit{test case}
          berhasil lulus tanpa kegagalan, mengindikasikan bahwa implementasi
          sistem manajemen perpustakaan telah memenuhi spesifikasi yang
          ditetapkan dan logika kondisionalnya berfungsi dengan benar.

    \item \textbf{Penerapan Prinsip \textit{Dependency Injection}.}
          Desain \texttt{LibraryControl} dengan \textit{constructor injection}
          memfasilitasi \textit{testability} kode, karena dependensi dapat
          dengan mudah digantikan oleh \textit{mock object} selama pengujian
          tanpa modifikasi kode produksi.
\end{enumerate}

% ─────────────────────────────────────────────────────────────────────────────
\section*{Referensi}
\addcontentsline{toc}{section}{Referensi}
% ─────────────────────────────────────────────────────────────────────────────

\begin{enumerate}[label={[\arabic*]}]
    \item Mockito Framework. \textit{Mockito -- Tasty mocking framework for
          unit tests in Java}. Tersedia di: \url{https://site.mockito.org/}
          [Diakses: Maret 2026].

    \item JUnit Team. \textit{JUnit 5 User Guide}. Tersedia di:
          \url{https://junit.org/junit5/docs/current/user-guide/}
          [Diakses: Maret 2026].

    \item Meszaros, G. (2007). \textit{xUnit Test Patterns: Refactoring Test
          Code}. Addison-Wesley Professional.

    \item Freeman, S., \& Pryce, N. (2009). \textit{Growing Object-Oriented
          Software, Guided by Tests}. Addison-Wesley Professional.

    \item Osherove, R. (2013). \textit{The Art of Unit Testing} (2nd ed.).
          Manning Publications.

    \item Apache Maven Project. \textit{Maven -- Welcome to Apache Maven}.
          Tersedia di: \url{https://maven.apache.org/}
          [Diakses: Maret 2026].
\end{enumerate}

\end{document}
